\MathBookSourcePage{102}
接着求解齐次边值问题
\[
\left\{
\begin{aligned}
-(\nu\Delta+r\operatorname{grad}\operatorname{div})
\mathbf{z}^m
&=\operatorname{grad}\omega^m
&&\text{in }\Omega,\\
\mathbf{z}^m&=\mathbf{0}
&&\text{on }\Gamma,
\end{aligned}
\right.
\]
并由
\[
\begin{aligned}
\rho^m
&=-\frac{\|\operatorname{div}\mathbf{u}^m\|_{0,\Omega}^2}
{(\operatorname{div}\mathbf{z}^m,
\operatorname{div}\mathbf{u}^m)},\\
p^{m+1}&=p^m-\rho^m\omega^m,\\
\mathbf{u}^{m+1}&=\mathbf{u}^m+\rho^m\mathbf{z}^m
\end{aligned}
\]
计算 $\rho^m\in\mathbb{R}$ 及下一对
$(\mathbf{u}^{m+1},p^{m+1})
\in H^1(\Omega)^N\times L_0^2(\Omega)$.
这里 Corollary \ref{cor:abstract-conjugate-gradient-convergence}
再次保证共轭梯度算法收敛.

本小节最后给出 Cattabriga [18] 的一个定理.
当边界 $\Gamma$ 足够光滑时, 该定理讨论 Stokes 问题
\eqref{eq:stokes-dirichlet-problem} 在更一般 Sobolev 空间中解的
存在性和正则性.

\setcounter{statement}{3}
\begin{Theorem}\label{thm:stokes-cattabriga-regularity}
除 Theorem \ref{thm:stokes-dirichlet-well-posedness} 的假设外,
再假设边界 $\Gamma$ 属于
$\mathcal{C}^{\max(2,m+2)}$ 类,
$\mathbf{f}\in W^{m,r}(\Omega)^N$ 且
$\mathbf{g}\in W^{m+2-1/r,r}(\Gamma)^N$,
其中整数 $m\ge-1$, 实数 $r$ 满足 $1<r<\infty$.
则问题 \eqref{eq:stokes-dirichlet-problem} 有唯一解
\[
(\mathbf{u},p)\in
W^{m+2,r}(\Omega)^N\times W^{m+1,r}(\Omega),
\qquad
\int_\Omega p\,\mathrm{d}x=0,
\]
并且存在与 $\mathbf{f},\mathbf{g}$ 无关的常数 $C$, 使
\begin{equation}\label{eq:stokes-cattabriga-regularity-bound}
\|\mathbf{u}\|_{m+2,r,\Omega}
+\|p\|_{m+1,r,\Omega}
\le
C\left(
\|\mathbf{f}\|_{m,r,\Omega}
+\|\mathbf{g}\|_{m+2-1/r,r,\Gamma}
\right).
\end{equation}
\end{Theorem}

\setcounter{statement}{5}
\begin{Remark}\label{rem:stokes-polyhedral-regularity}
当 $\Omega$ 是多面体区域时, 问题
\eqref{eq:stokes-dirichlet-problem} 的解具有明显较弱的正则性.
特别地, 当 $N=2$ 且 $\Omega$ 是凸多边形时,
对 $m\le0$ 和 $1<r\le2$, Theorem
\ref{thm:stokes-cattabriga-regularity} 的结论仍然成立
(参见 Grisvard [43]).
\end{Remark}

\subsection{二维 Dirichlet 问题的流函数形式}
\label{subsec:stokes-dirichlet-stream-function}

仍像 Figure \ref{fig:sec2-normal-tangent} 那样, 用
$\Gamma_i$, $0\le i\le p$, 表示边界 $\Gamma$ 的各连通分支.
现在不用 \eqref{eq:stokes-boundary-flux-compatibility},
而假设更强的条件
\begin{equation}\label{eq:stokes-componentwise-flux-compatibility}
\int_{\Gamma_i}\mathbf{g}\cdot\mathbf{n}\,\mathrm{d}s=0,
\qquad 0\le i\le p.
\end{equation}
于是根据 Theorem \ref{thm:sec3-stream-function-characterization}
和 Theorem \ref{thm:sec3-three-dimensional-vector-potential},
Theorem \ref{thm:stokes-dirichlet-well-posedness} 给出的速度场
$\mathbf{u}$ 可表示为流函数 $\psi$ 的 curl($N=2$),
或向量势 $\boldsymbol{\psi}$ 的 curl($N=3$).
下面说明这个流函数或向量势可刻画为 $\Omega$ 中一个双调和问题的解.

先考虑 $N=2$. 流函数 $\psi$ 在相差一个加性常数的意义下唯一.
但由于
$\psi\in H^2(\Omega)\subset\mathcal{C}^0(\overline{\Omega})$,
可以通过固定 $\psi$ 在 $\overline{\Omega}$ 中某一点的值来唯一确定它.
首先置 $\psi(x_0)=0$, 其中 $x_0$ 是 $\Gamma_0$ 上的任意一点
(例如!). 然后选取满足
\begin{equation}\label{eq:stokes-stream-boundary-primitive}
\frac{\partial\chi}{\partial\tau}
=\mathbf{g}\cdot\mathbf{n}
\quad\text{on }\Gamma,
\qquad
\chi(x_0)=0
\end{equation}
的函数 $\chi\in H^{3/2}(\Gamma)$.

\MathBookSourcePage{103}
由于
$\partial\psi/\partial\tau
=\mathbf{g}\cdot\mathbf{n}$ 在 $\Gamma$ 上成立, 可知
\[
\psi=
\left\{
\begin{aligned}
\chi&&\text{on }\Gamma_0,\\
\chi+c_i&&\text{on }\Gamma_i,
\quad 1\le i\le p,
\end{aligned}
\right.
\]
其中常数 $c_i$ 是确定但未知的.

\setcounter{statement}{4}
\begin{Theorem}\label{thm:stokes-two-dimensional-stream-function}
设 $N=2$, 并且 Theorem
\ref{thm:stokes-dirichlet-well-posedness} 的假设连同
\eqref{eq:stokes-componentwise-flux-compatibility} 一起成立.
则 $\mathbf{u}$ 的相应流函数可刻画为下列方程的唯一解
$\psi\in H^2(\Omega)$:
\begin{equation}\label{eq:stokes-stream-biharmonic-variational}
\nu(\Delta\psi,\Delta\phi)
=\langle\mathbf{f},\operatorname{curl}\phi\rangle
\qquad
\forall\phi\in\Psi
\quad
(\text{cf. }\eqref{eq:sec3-stream-variational}),
\end{equation}
\begin{equation}\label{eq:stokes-stream-boundary-conditions}
\left\{
\begin{aligned}
\psi&=\chi
&&\text{on }\Gamma_0,\\
\psi&=\chi+c_i
&&\text{on }\Gamma_i,
\quad 1\le i\le p,\\
\frac{\partial\psi}{\partial n}
&=-\mathbf{g}\cdot\boldsymbol{\tau}
&&\text{on }\Gamma,
\end{aligned}
\right.
\end{equation}
其中 $\chi$ 按照
\eqref{eq:stokes-stream-boundary-primitive} 选取.
\end{Theorem}
\begin{Proof}
已经证明速度场 $\mathbf{u}$ 是
\eqref{eq:stokes-dirichlet-p-formulation} 的唯一解.
函数 $\mathbf{u}$ 满足
\[
\operatorname{div}\mathbf{u}=0\quad\text{in }\Omega,
\qquad
\mathbf{u}=\mathbf{g}\quad\text{on }\Gamma
\]
当且仅当 $\mathbf{u}=\operatorname{curl}\psi$,
其中流函数 $\psi\in H^2(\Omega)$ 满足边界条件
\eqref{eq:stokes-stream-boundary-conditions}.
此外, 根据 Corollary \ref{cor:sec3-v-stream-functions},
$\mathbf{v}\in V$ 当且仅当
$\mathbf{v}=\operatorname{curl}\phi$ 且 $\phi\in\Psi$.
所以只需证明
\begin{equation}\label{eq:stokes-stream-gradient-laplacian-identity}
(\operatorname{grad}\mathbf{u},
\operatorname{grad}\mathbf{v})
=
(\Delta\psi,\Delta\phi)
\qquad
\forall\mathbf{v}=\operatorname{curl}\phi,
\quad \phi\in\Psi.
\end{equation}
这里使用 Subsection \ref{subsec:sec2-curl-spaces} 开头给出的恒等式.
首先
\[
\begin{aligned}
(\Delta\psi,\Delta\phi)
&=(\operatorname{curl}(\operatorname{curl}\psi),
\operatorname{curl}(\operatorname{curl}\phi))\\
&=(\operatorname{curl}\mathbf{u},
\operatorname{curl}\mathbf{v}).
\end{aligned}
\]
其次取
\[
\mathbf{v}\in\mathcal{V}
=\{\mathbf{v}\in\mathcal{D}(\Omega)^2;
\ \operatorname{div}\mathbf{v}=0\},
\]
并回忆 Corollary \ref{cor:sec2-test-divfree-dense} 说明
$\mathcal{V}$ 在 $V$ 中稠密. 在分布意义下,
\[
\begin{aligned}
(\operatorname{curl}\mathbf{u},
\operatorname{curl}\mathbf{v})
&=\langle
\operatorname{curl}(\operatorname{curl}\mathbf{u}),
\mathbf{v}\rangle\\
&=\langle-\Delta\mathbf{u},\mathbf{v}\rangle
\quad\text{because }\operatorname{div}\mathbf{u}=0\\
&=(\operatorname{grad}\mathbf{u},
\operatorname{grad}\mathbf{v}).
\end{aligned}
\]
因此 $\mathcal{V}$ 在 $V$ 中的稠密性给出
\eqref{eq:stokes-stream-gradient-laplacian-identity}.
\end{Proof}

还需要解释问题
\eqref{eq:stokes-stream-biharmonic-variational}--
\eqref{eq:stokes-stream-boundary-conditions}.
形式地应用 Green 公式, 容易说明 $\psi$ 是边值问题的唯一解:
\[
\nu\Delta^2\psi=\operatorname{curl}\mathbf{f},
\]
\[
\psi=\chi\quad\text{on }\Gamma_0,
\qquad
\psi=c_i+\chi\quad\text{on }\Gamma_i,
\quad 1\le i\le p,
\]
\[
\frac{\partial\psi}{\partial n}
=-\mathbf{g}\cdot\boldsymbol{\tau}
\quad\text{on }\Gamma,
\]
\[
\int_{\Gamma_i}
\left(
\nu\frac{\partial(\Delta\psi)}{\partial n}
-\mathbf{f}\cdot\boldsymbol{\tau}
\right)\,\mathrm{d}s
=0,
\qquad 1\le i\le p.
\]
注意, 当 $\mathbf{f}\in H(\operatorname{curl};\Omega)$ 时,
最后一个条件有意义.

\MathBookSourcePage{104}
\subsection{三维情形}
\label{subsec:stokes-three-dimensional-vector-potential}

现在考虑 $N=3$. 为简化讨论, 先考察齐次边界条件
$\mathbf{g}=\mathbf{0}$ 的情形. 从 Subsection
\ref{subsec:three-dimensional-decomposition} 已知, 速度场
$\mathbf{u}$ 可以表示为
\[
\mathbf{u}=\operatorname{curl}\boldsymbol{\psi},
\qquad
\operatorname{div}\boldsymbol{\psi}=0
\quad\text{in }\Omega.
\]
当 $\Omega$ 单连通时, 向量势 $\boldsymbol{\psi}$ 可由下列两组条件
中的任一组唯一确定:
\begin{equation}\label{eq:stokes-vector-potential-normal-boundary}
\boldsymbol{\psi}\cdot\mathbf{n}=0
\quad\text{on }\Gamma,
\end{equation}
或者
\begin{equation}\label{eq:stokes-vector-potential-tangential-boundary}
\boldsymbol{\psi}\times\mathbf{n}=\mathbf{0}
\quad\text{on }\Gamma,
\qquad
\int_{\Gamma_i}\boldsymbol{\psi}\cdot\mathbf{n}\,\mathrm{d}s=0,
\quad 0\le i\le p.
\end{equation}
由于 $\mathbf{u}$ 在 $\Gamma$ 上为零, 可以添加边界条件
\begin{equation}\label{eq:stokes-vector-potential-curl-boundary}
\operatorname{curl}\boldsymbol{\psi}=\mathbf{0}
\quad\text{on }\Gamma,
\end{equation}
但注意, 当 $\boldsymbol{\psi}$ 满足
\eqref{eq:stokes-vector-potential-tangential-boundary} 的第一个条件时,
这一条件化为
\begin{equation}\tag{\ref*{eq:stokes-vector-potential-curl-boundary}$'$}
\label{eq:stokes-vector-potential-curl-tangential-boundary}
(\operatorname{curl}\boldsymbol{\psi})\times\mathbf{n}
=\mathbf{0}
\quad\text{on }\Gamma
\end{equation}
(参见 Remark \ref{rem:sec2-curl-in-h}).

关于 $\boldsymbol{\psi}$ 的正则性, 我们所需的只是
$\boldsymbol{\psi}$ 满足一个双调和问题. 因而要求
$\Delta\boldsymbol{\psi}\in L^2(\Omega)^3$ 看来是合理的.
由恒等式
\begin{equation}\label{eq:stokes-vector-laplacian-decomposition}
-\Delta\boldsymbol{\phi}
=
\operatorname{curl}\operatorname{curl}\boldsymbol{\phi}
-\operatorname{grad}(\operatorname{div}\boldsymbol{\phi})
\end{equation}
以及
$\operatorname{curl}\boldsymbol{\psi}\in H^1(\Omega)^3$,
可见这相当于要求
$\operatorname{div}\boldsymbol{\psi}\in H^1(\Omega)$.
当然, 当 $\boldsymbol{\psi}$ 无散时总是如此,
但在实际中无散条件并不方便, 下面将完全放松这一条件.

此时必须指出, 在多连通区域 $\Omega$ 上,
向量势 $\boldsymbol{\psi}$ 没有直接的刻画.
为此, 对后一种情形只作简要讨论, 并主要限于单连通区域.

先考虑满足
\eqref{eq:stokes-vector-potential-tangential-boundary} 的向量势.
根据上述考虑, 引入空间
\[
\begin{aligned}
\Psi=\{\,\boldsymbol{\phi}\in L^2(\Omega)^3;\;&
\operatorname{div}\boldsymbol{\phi}\in H^1(\Omega),\
\operatorname{curl}\boldsymbol{\phi}\in H_0^1(\Omega)^3,\\
&\boldsymbol{\phi}\times\mathbf{n}|_\Gamma=\mathbf{0},\
\int_{\Gamma_i}\boldsymbol{\phi}\cdot\mathbf{n}\,\mathrm{d}s=0,\
0\le i\le p\,\},
\end{aligned}
\]
其范数为
\[
\|\boldsymbol{\phi}\|
=
\left\{
\|\boldsymbol{\phi}\|_{0,\Omega}^2
+\|\operatorname{div}\boldsymbol{\phi}\|_{1,\Omega}^2
+\|\operatorname{curl}\boldsymbol{\phi}\|_{1,\Omega}^2
\right\}^{1/2}.
\]

\MathBookSourcePage{105}
注意, $\Psi$ 中的函数并非无散, 但满足
$\int_\Omega\operatorname{div}\boldsymbol{\phi}\,\mathrm{d}x=0$.

现在设 $\mathbf{u}$ 是边界条件 $\mathbf{g}=\mathbf{0}$ 时问题
\eqref{eq:stokes-dirichlet-problem} 的解. 如上所述,
当 $\Omega$ 单连通时, $\mathbf{u}$ 在 $\Psi$ 中有唯一的无散向量势
$\boldsymbol{\psi}$:
\begin{equation}\label{eq:stokes-vector-potential-stokes-system}
\left\{
\begin{aligned}
-\nu\Delta(\operatorname{curl}\boldsymbol{\psi})
+\operatorname{grad}p
&=\mathbf{f}
&&\text{in }H^{-1}(\Omega)^3,\\
\operatorname{div}\boldsymbol{\psi}&=0
&&\text{in }\Omega,\qquad
\boldsymbol{\psi}\in\Psi.
\end{aligned}
\right.
\end{equation}
因为
$\operatorname{div}(\operatorname{curl}\boldsymbol{\psi})=0$,
\eqref{eq:stokes-vector-potential-stokes-system} 化为
\[
\nu\operatorname{curl}\operatorname{curl}
(\operatorname{curl}\boldsymbol{\psi})
+\operatorname{grad}p=\mathbf{f}.
\]
将方程两边与
$\operatorname{curl}\boldsymbol{\phi}$ 相乘,
其中 $\boldsymbol{\phi}\in\Psi$. 由于
$\operatorname{curl}\boldsymbol{\phi}\in H_0^1(\Omega)^3$, 有
\[
\begin{aligned}
\langle
\operatorname{curl}\operatorname{curl}
(\operatorname{curl}\boldsymbol{\psi}),
\operatorname{curl}\boldsymbol{\phi}
\rangle
=
(&\operatorname{curl}(\operatorname{curl}\boldsymbol{\psi}),\\
&\operatorname{curl}(\operatorname{curl}\boldsymbol{\phi})).
\end{aligned}
\]
所以只需证明
\begin{equation}\label{eq:stokes-vector-curlcurl-graddiv-orthogonality}
(\operatorname{curl}\operatorname{curl}\boldsymbol{\psi},
\operatorname{grad}\operatorname{div}\boldsymbol{\phi})=0
\end{equation}
便可得到
\[
-\langle
\Delta(\operatorname{curl}\boldsymbol{\psi}),
\operatorname{curl}\boldsymbol{\phi}
\rangle
=
(\Delta\boldsymbol{\psi},
\Delta\boldsymbol{\phi})
\qquad
\forall\boldsymbol{\phi}\in\Psi.
\]
只要
$\operatorname{curl}\boldsymbol{\psi}\in H_0^1(\Omega)^3$,
\eqref{eq:stokes-vector-curlcurl-graddiv-orthogonality} 即成立.
于是已经证明 $\boldsymbol{\psi}$ 是双调和问题

求 $\boldsymbol{\psi}\in\Psi$, 使
\begin{equation}\label{eq:stokes-vector-biharmonic-tangential}
\nu(\Delta\boldsymbol{\psi},
\Delta\boldsymbol{\phi})
=
\langle\mathbf{f},
\operatorname{curl}\boldsymbol{\phi}\rangle
\qquad
\forall\boldsymbol{\phi}\in\Psi
\end{equation}
的解.

反过来, 容易证明
\eqref{eq:stokes-vector-biharmonic-tangential} 至多有一个解.
事实上, 若 $\boldsymbol{\psi}\in\Psi$ 满足
$\Delta\boldsymbol{\psi}=\mathbf{0}$,
则 $\operatorname{curl}\boldsymbol{\psi}=\mathbf{0}$ in $\Omega$,
因为 $\operatorname{curl}\boldsymbol{\psi}$ 是 Dirichlet 问题
\[
\Delta(\operatorname{curl}\boldsymbol{\psi})=\mathbf{0}
\quad\text{in }\Omega,
\qquad
\operatorname{curl}\boldsymbol{\psi}|_\Gamma=\mathbf{0}
\]
的解. 进而
$\operatorname{grad}(\operatorname{div}\boldsymbol{\psi})=\mathbf{0}$.
由于
$\int_\Omega\operatorname{div}\boldsymbol{\psi}\,\mathrm{d}x=0$,
所以 $\operatorname{div}\boldsymbol{\psi}=0$.
因此 $\boldsymbol{\psi}=\mathbf{0}$(参见 Remark
\ref{rem:sec3-zero-curl-zero-tangential-implies-zero}). 由此证明下述结果.

\setcounter{statement}{0}
\begin{Lemma}\label{lem:stokes-vector-biharmonic-tangential}
若 $\Omega$ 单连通并满足 Theorem
\ref{thm:stokes-dirichlet-well-posedness} 的假设, 则双调和问题
\eqref{eq:stokes-vector-biharmonic-tangential} 有唯一解
$\boldsymbol{\psi}$. 此外,
$\operatorname{div}\boldsymbol{\psi}=0$, 且
$\operatorname{curl}\boldsymbol{\psi}$ 是边界条件
$\mathbf{g}=\mathbf{0}$ 时 Stokes 问题
\eqref{eq:stokes-dirichlet-problem} 的唯一解.
\end{Lemma}
